\documentclass[12pt]{article}
\usepackage{amsmath, amssymb, amsthm}
\usepackage{mathtools}
\usepackage{tensor}
\usepackage{geometry}
\usepackage{booktabs}
\usepackage{array}
\usepackage{xcolor}
\usepackage{hyperref}
\usepackage{fancyhdr}
\usepackage{slashed}
\usepackage{mdframed}
\geometry{margin=1.1in, headheight=15pt}

\definecolor{cadblue}{RGB}{30,90,200}
\definecolor{verifygreen}{RGB}{20,140,60}

\newcommand{\cdbexpr}[1]{\textcolor{cadblue}{$#1$}}
\newcommand{\verified}[1]{\textcolor{verifygreen}{\checkmark\; #1}}
\newcommand{\half}{\tfrac{1}{2}}
\DeclareMathOperator\Tr{Tr}
\newtheorem{theorem}{Theorem}

\pagestyle{fancy}
\fancyhf{}
\lhead{Srednicki QFT --- Chapter 43}
\rhead{Feynman Rules for Majorana Fields}
\cfoot{\thepage}

\title{\textbf{Srednicki QFT: Chapter 43}\\[6pt]
       \large The Feynman Rules for Majorana Fields\\[4pt]
       \normalsize Majorana condition, propagator, Feynman rules,\\
       and the connection to MHV amplitudes in SUSY QCD}
\author{Cadabra2 + NumPy verification}
\date{}

\begin{document}
\maketitle
\tableofcontents
\newpage

\begin{mdframed}[backgroundcolor=blue!5]
\textbf{Chapter 43} covers a special spinor: the \emph{Majorana field},
where the particle is its own antiparticle.  Majorana fermions appear in:
supersymmetric theories (gauginos, neutralinos), potential neutrino mass terms,
and as a theoretical tool in the construction of MHV amplitudes in $\mathcal{N}=4$ SYM.
\end{mdframed}

%% ============================================================
\section{The Majorana Condition}
%% ============================================================

\textit{Srednicki eq.\ (43.1)--(43.5).}

\begin{mdframed}[backgroundcolor=blue!5]
\textbf{Majorana condition:}
\begin{equation}\label{eq:majorana}
  \psi(x) = \psi^c(x) \equiv C\,\bar\psi^T(x),
\end{equation}
where $C = i\gamma^2\gamma^0$ is the charge-conjugation matrix satisfying:
\begin{align}
  C^2 &= -\mathbf{1}, \\
  C\gamma^\mu C^{-1} &= -(\gamma^\mu)^T, \\
  C^T &= -C, \quad C^\dagger = C^{-1}.
\end{align}
\end{mdframed}

The Majorana condition \eqref{eq:majorana} states that the particle is
\emph{its own antiparticle}: there is no distinction between $\psi$ and $\psi^c$.
This halves the degrees of freedom from 4 (Dirac) to 2.

\subsection{Comparison of spinor types}

\begin{center}
\begin{tabular}{lccc}
\toprule
Property & Weyl & Majorana & Dirac \\
\midrule
Degrees of freedom & 2 & 2 & 4 \\
EM charge & 0 & 0 & $\pm e$ \\
Particle $=$ antiparticle & --- & \checkmark & $\times$ \\
Lorentz rep. & $(\half,0)$ or $(0,\half)$ & $(0,\half)\oplus(\half,0)$ & $(0,\half)\oplus(\half,0)$ \\
SUSY partner & scalar & gaugino & sfermion pair \\
\bottomrule
\end{tabular}
\end{center}

\subsection{Majorana from Weyl}

A Majorana spinor can be constructed from a single left-handed Weyl spinor $\chi_\alpha$:
\begin{equation}
  \psi_M = \begin{pmatrix}\chi_\alpha \\ \bar\chi^{\dot\alpha}\end{pmatrix}
  = \begin{pmatrix}\chi \\ i\sigma^2\chi^*\end{pmatrix}.
\end{equation}
This is automatically self-conjugate: $\psi_M = C\bar\psi_M^T$.

%% ============================================================
\section{Majorana Lagrangian and Mass}
%% ============================================================

\textit{Srednicki eq.\ (43.6)--(43.11).}

The free Majorana Lagrangian:
\begin{equation}
  \mathcal{L}_M = \frac{i}{2}\psi^T C\slashed\partial\psi
               - \frac{m_M}{2}\psi^T C\psi + \text{h.c.}
\end{equation}

\begin{mdframed}[backgroundcolor=blue!5]
\textbf{Majorana mass term:}
\begin{equation}\label{eq:majorana-mass}
  \mathcal{L}_\text{mass} = -\frac{m_M}{2}\psi^T C\psi + \text{h.c.}
  = -\frac{m_M}{2}(\chi\chi + \bar\chi\bar\chi),
\end{equation}
where the $\tfrac{1}{2}$ reflects particle $=$ antiparticle (no double-counting).
\end{mdframed}

MSSM examples:
\begin{align}
  \text{Gaugino masses:} &\quad \frac{1}{2}M_1\tilde\lambda_B\tilde\lambda_B
                                + \frac{1}{2}M_2\tilde\lambda_W\tilde\lambda_W
                                + \frac{1}{2}M_3\tilde g\tilde g + \text{h.c.},\\
  \text{Higgsino masses:} &\quad \mu H_u H_d + \text{h.c.}
\end{align}

%% ============================================================
\section{The Majorana Propagator}
%% ============================================================

\textit{Srednicki eq.\ (43.12)--(43.15).}

\begin{mdframed}[backgroundcolor=blue!5]
\textbf{Majorana propagator:}
\begin{equation}\label{eq:majorana-prop}
  S_F^M(k) = \frac{i(\slashed{k}+m)}{k^2+m^2-i\varepsilon}.
\end{equation}
The Wick contraction produces a factor of $\tfrac{1}{2}$ relative to Dirac:
\begin{equation}
  \langle T\,\psi_\alpha(x)\bar\psi_\beta(y)\rangle_M
  = \tfrac{1}{2}\,\langle T\,\psi_\alpha(x)\bar\psi_\beta(y)\rangle_D,
\end{equation}
because each fermion line can be traversed in only one direction (particle $=$ antiparticle).
\end{mdframed}

%% ============================================================
\section{Feynman Rules for Majorana Fields}
%% ============================================================

\textit{Srednicki eq.\ (43.16)--(43.22).}

\subsection{Standard rules}

\begin{center}
\begin{tabular}{lc}
\toprule
Object & Feynman rule \\
\midrule
Majorana propagator & $\displaystyle\frac{i(\slashed{k}+m)}{k^2+m^2-i\varepsilon}$ \\[8pt]
External incoming Majorana & $u_s(\mathbf{p})$ or $v_s(\mathbf{p})$ \\
External outgoing Majorana & $\bar u_s(\mathbf{p})$ or $\bar v_s(\mathbf{p})$ \\
Yukawa vertex $g\phi\psi\psi$ & $-ig$ \\
4-Majorana contact $\lambda(\psi\psi)^2$ & $-4i\lambda$ \\
Closed Majorana loop & $+\tfrac{1}{2}\Tr[\cdots]$ \\
\bottomrule
\end{tabular}
\end{center}

\subsection{Sign rules}

Majorana spinors obey the same anticommutation statistics as Dirac.  For external
legs the relative sign between diagrams is determined by the fermion-flow arrows:
when two diagrams differ by an exchange of identical Majorana particles, a relative
$(-1)$ is required.

%% ============================================================
\section{Why Majorana Matters for MHV}
%% ============================================================

\textit{Connecting Ch.~43 to Ch.~48 and beyond.}

\subsection{Gluinos are Majorana}

In $\mathcal{N}=1$ Super QCD (SQCD), the gluino $\tilde g^a$ is the superpartner
of the gluon.  It transforms in the adjoint of $SU(N_c)$ and is \emph{Majorana}:
\begin{equation}
  \tilde g^a = C\,\overline{\tilde g}^{\,a\,T}.
\end{equation}
This follows from the gauge symmetry of the $\mathcal{N}=1$ vector multiplet:
since the gauge field $A_\mu^a$ is real (no particle/antiparticle distinction in
the adjoint representation), its superpartner must also be self-conjugate.

\subsection{Spin sums for Majorana spinors}

\begin{mdframed}[backgroundcolor=blue!5]
\textbf{Key:} For a Majorana fermion, particle and antiparticle states are
identical, so the spin sum is:
\begin{equation}\label{eq:majorana-spinsum}
  \sum_s u_s(\mathbf{p})\bar u_s(\mathbf{p}) = -\slashed{p} + m,
\end{equation}
\emph{exactly} as for a Dirac fermion.  The Majorana condition does not modify
the spin-sum completeness relation.
\end{mdframed}

The proof is straightforward: since $v_s(\mathbf{p}) = C\bar u_s^T(\mathbf{p})$
for a Majorana spinor, the $v$-sum equals the $u$-sum, and both give $-\slashed{p}+m$
(massless: $\sum_s u_s\bar u_s = -\slashed{p}$).

\subsection{The MHV amplitude with a gluino pair}

Consider the MHV amplitude in $\mathcal{N}=4$ SYM with two external gluinos
(Majorana, adjoint) at positions $j,k$ and $n-2$ positive-helicity gluons:

\begin{mdframed}[backgroundcolor=blue!5]
\begin{equation}\label{eq:mhv-gluino}
  A_n^{\rm MHV}(1^+,\ldots,j^{\tilde g-},\ldots,k^{\tilde g+},\ldots,n^+)
  = i\,\frac{\langle jk\rangle^3 [jk]}
            {\langle12\rangle\langle23\rangle\cdots\langle n1\rangle}.
\end{equation}
\end{mdframed}

Comparing with the pure-gluon MHV formula $A_n^{\rm MHV} = i\langle jk\rangle^4/(\ldots)$,
the gluino numerator is $\langle jk\rangle^3[jk]$ instead of $\langle jk\rangle^4$.
This is because:
\begin{enumerate}
  \item Each negative-helicity gluon contributes $\langle jk\rangle^2$ via
        $\Tr[\slashed{p}_j(1-\gamma^5)\slashed{p}_k(1+\gamma^5)] = 4\langle jk\rangle^2$
        from the 4-gamma trace identity (Ch.~41).
  \item A negative-helicity \emph{Majorana} gluino contributes $\langle jk\rangle$
        (one power less) because it carries helicity $h=-\tfrac12$ instead of $h=-1$.
  \item The $[jk]$ factor comes from the spin sum \eqref{eq:majorana-spinsum}: since
        $\sum_s u_s\bar u_s = -\slashed{p}$ and $p_{\alpha\dot\alpha} = \lambda_\alpha\tilde\lambda_{\dot\alpha}$,
        the $\tilde\lambda$ factor of leg $j$ combined with the $\lambda$ factor of leg $k$
        gives $[jk] = \tilde\lambda_j^{\dot\alpha}\tilde\lambda_{k\dot\alpha}$.
\end{enumerate}

\subsection{The Majorana $\tfrac{1}{2}$ factor and loop amplitudes}

For loop amplitudes in SQCD/SYM, Majorana fermion loops contribute with a factor of
$\tfrac{1}{2}$ relative to Dirac fermions.  This is essential for:
\begin{itemize}
  \item The one-loop beta function in $\mathcal{N}=1$ SQCD:
        $\beta(g) = -\frac{g^3}{16\pi^2}(3N_c - N_f)$, where the $3N_c$ from gluino
        loops uses the Majorana $\tfrac{1}{2}$.
  \item Unitarity cuts: a Majorana fermion in the cut contributes $\tfrac{1}{2}\Tr[\slashed{l}]$
        compared to a Dirac fermion's $\Tr[\slashed{l}]$.
\end{itemize}

\subsection{BCFW for Majorana}

The Britto-Cachazo-Feng-Witten (BCFW) recursion relation extends naturally to Majorana
fermions.  For a shift $\hat\lambda_j \to \lambda_j + z\lambda_k$,
$\hat{\tilde\lambda}_k \to \tilde\lambda_k - z\tilde\lambda_j$ (a $[jk\rangle$ shift):
\begin{equation}
  A_n = \sum_{\text{diagrams}} A_L(\hat j, P) \frac{i}{P^2} A_R(-P, \hat k),
\end{equation}
where the sum over internal helicities for a Majorana intermediate state
uses $\sum_s u_s\bar u_s = -\slashed{P}$ — \emph{the same completeness relation
as for Dirac}, so the BCFW recursion is identical in form.

\begin{mdframed}[backgroundcolor=blue!5]
\textbf{Summary: Majorana $\leftrightarrow$ MHV connections}
\begin{center}
\begin{tabular}{ll}
\toprule
Majorana property & Role in MHV/BCFW \\
\midrule
$\psi = C\bar\psi^T$ (self-conjugate) & No distinction $u/v$ in spin sum \\
$\sum_s u_s\bar u_s = -\slashed{p}$ & Same $\slashed{p}$ as in Parke-Taylor spinors \\
Wick $\tfrac{1}{2}$ factor & $\tfrac{1}{2}\Tr[\cdots]$ for Majorana loops \\
Gluino in $\mathcal{N}=4$ SYM & $\langle jk\rangle^3[jk]$ vs.\ $\langle jk\rangle^4$ for gluons \\
BCFW recursion & Identical form; Majorana = Dirac at tree level \\
\bottomrule
\end{tabular}
\end{center}
\end{mdframed}

\vfill
\bibliographystyle{plain}
\end{document}
